\documentclass[10pt,a4paper]{amsart}
\usepackage[margin=19mm]{geometry}
\usepackage{amsmath,amssymb,mathtools}
\usepackage[scr=rsfs,cal=euler]{mathalfa}
\usepackage{xfrac}
\usepackage[unicode,hidelinks,bookmarks=false]{hyperref}
\newcommand{\ie}{i.e.\ }
\newcommand{\cf}{cf.\ }
\newcommand{\resp}{respectively\ }
\newcommand{\Subsection}{\subsection}
\providecommand{\nobreakdash}{\nobreak\hbox{-}}
\setlength{\emergencystretch}{2em}
\multlinegap0pt
\usepackage{amssymb}
\usepackage[shortlabels]{enumitem}
\usepackage{tikz}
\usepackage{algorithm}\usepackage{algpseudocode}
\newtheorem{lemma}{Lemma}
\newtheorem{theorem}{Theorem}
\newcommand{\TV}{\mathrm{TV}}
\title{Residual Centering and Error Bounds for Convex Approximation of Mixed-Integer Recourse}
\author{Alban Kryeziu}
\date{Source: arXiv:2608.00245v1 (CC BY 4.0)}
\begin{document}
\maketitle

\noindent\emph{Source and licence.} Residual Centering and Error Bounds for Convex Approximation of Mixed-Integer Recourse. Version arXiv:2608.00245v1 (CC BY 4.0), distributed by arXiv under the Creative Commons Attribution 4.0 International licence, http://creativecommons.org/licenses/by/4.0/, as recorded in the arXiv licence metadata for this exact version and shown on the abstract page https://arxiv.org/abs/2608.00245v1. Full licence notice: \texttt{licenses/arxiv-2608.00245-CC-BY-4.0.txt}.\par\smallskip

\noindent\textit{Editorial scope.} Counted unit: the article's theorem thm:k\_direction\_tv\_bound (source0.tex lines 1088--1119). The article's complete original proof of it is delivered with it (source0.tex lines 1120--1293, one \texttt{proof} environment of 174 lines). No mathematics is written, repaired or re-derived: the statement and the proof are the article's own lines, in the article's own notation, and no retained line is substituted or re-flowed.  Retained uncounted as the source-local context the proof needs: lines 1026--1043.  Nothing was omitted from any retained span, so the package carries no omission record.  No retained line contains a citation, so the wrapper carries no bibliography and no bibliography entry.  No retained line contains a figure environment, an image-inclusion command or drawing code, so no image is delivered and the package declares no asset dependency.  Editorial formatting only, in addition: an AMS \texttt{amsart} preamble that restates the article's own declarations and macros used by the retained lines, namely the environments \texttt{lemma}, \texttt{theorem} on separate counters, together with the standard packages those lines need.  The retained environments print in this document as Lemma 1, Theorem 1 in order of appearance, whereas the article prints the same items with the numbering of its own full document.  The complete native source of the article as distributed in the arXiv e-print for this exact version is bundled unchanged as \texttt{sources/206528/source0.tex}.
\begin{lemma}[Half-oscillation inequality]
\label{lem:half_oscillation}
Let $g\in L^1([a,b])$ satisfy $\int_a^b g(t)\,dt=0$, and define
$G(t):=\int_a^t g(s)\,ds$. Then
\begin{equation}
\label{eq:half_oscillation}
\|G\|_{L^\infty([a,b])}
\le
\frac12\int_a^b |g(s)|\,ds.
\end{equation}
Moreover,
\begin{equation}
\label{eq:recentered_half_oscillation}
\inf_{c\in\mathbb R}\|G-c\|_{L^\infty([a,b])}
\le
\frac14\int_a^b |g(s)|\,ds.
\end{equation}
\end{lemma}

\begin{theorem}[$k$-direction anisotropic $\TV_\infty$ bound]
\label{thm:k_direction_tv_bound}
Let $I\subseteq[d]$ be nonempty and set $k:=|I|$. Let $f\in BV(S)$ be a
probability density on $S$, and let $\varphi\in W^{1,\infty}(S)$ be centered on
$I$. Then
\begin{equation}
\label{eq:k_direction_tv_bound}
|\mathbb E_f[\varphi]|
\le
\frac{1}{2k}
\Lambda_{\mathrm{sl}}^I(\varphi)
\TV_I(f;S).
\end{equation}
Consequently,
\begin{equation}
\label{eq:k_direction_tv_bound_tvinfty}
|\mathbb E_f[\varphi]|
\le
\frac{1}{2k}
\Lambda_{\mathrm{sl}}^I(\varphi)
\TV_\infty(f;S).
\end{equation}
If $\|\varphi\|_{L^\infty(S)}\le h$, then
\begin{equation}
\label{eq:k_direction_tv_bound_sup}
|\mathbb E_f[\varphi]|
\le
\frac{h}{2k}
\left(\max_{i\in I}(b_i-a_i)\right)
\TV_I(f;S).
\end{equation}
\end{theorem}

\begin{proof}
We subtract the uniform density to remove the constant part of $f$. This does
not change the variation inside $S$, because $u$ has no distributional
derivative in the interior of the box. Let $g:=f-u$. Since $u$ is constant on
$S$, its distributional derivatives inside $S$ vanish. Hence
$$
D_i g=D_i f
\qquad\text{as measures on }S,
$$
for each coordinate $i$, and therefore
$$
\TV_I(g;S)=\TV_I(f;S).
$$

Because $I$ is nonempty and $\varphi$ is centered on $I$, choosing any
$i_0\in I$ and applying Fubini's theorem gives
$$
\int_S\varphi(\omega)\,d\omega=0.
$$
Thus $\mathbb E_u[\varphi]=0$, and hence
$$
\mathbb E_f[\varphi]
=
\mathbb E_f[\varphi]-\mathbb E_u[\varphi]
=
\int_S (f-u)\varphi\,d\omega
=
\int_S g\varphi\,d\omega.
$$

Define a vector field $\Psi:\overline S\to\mathbb R^d$ by
$$
\Psi_i(\omega)
:=
\frac1k
\int_{a_i}^{\omega_i}
\varphi(\omega_1,\ldots,\omega_{i-1},t,\omega_{i+1},\ldots,\omega_d)\,dt,
\qquad i\in I,
$$
and set $\Psi_i(\omega):=0$ for $i\notin I$. Since
$\varphi\in W^{1,\infty}(S)$, we use its Lipschitz representative on
$\overline S$. Then $\Psi\in\mathrm{Lip}(\overline S;\mathbb R^d)$. For
$i\in I$,
$$
\partial_i\Psi_i(\omega)=\frac1k\varphi(\omega)
\qquad\text{for a.e. }\omega\in S,
$$
and therefore
$$
\operatorname{div}\Psi
=
\sum_{i\in I}\partial_i\Psi_i
=
\varphi
\qquad\text{a.e. on }S.
$$

The key point is that the normal trace of $\Psi$ vanishes, and this is
exactly where the centering assumption is used. If $i\in I$, then $\Psi_i=0$ on the lower
face $\{\omega_i=a_i\}$. On the upper face $\{\omega_i=b_i\}$,
$$
\Psi_i(\omega_1,\ldots,b_i,\ldots,\omega_d)
=
\frac1k
\int_{a_i}^{b_i}
\varphi(\omega_1,\ldots,t,\ldots,\omega_d)\,dt,
$$
which is zero for a.e. $\omega_{-i}$ because $\varphi$ is centered in direction
$i$. If $j\notin I$, then $\Psi_j$ is identically zero. Since the outward
normal on the two faces orthogonal to $e_j$ is $\pm e_j$, it follows that
$$
\Psi\cdot n=0
\qquad
\mathcal H^{d-1}\text{-a.e. on }\partial S.
$$

It remains to estimate the size of $\Psi$. Fix $i\in I$ and a slice
$\omega_{-i}\in S_{-i}$ for which the centering identity holds. Set
$$
h_{\omega_{-i}}(t)
:=
\varphi(\omega_1,\ldots,\omega_{i-1},t,\omega_{i+1},\ldots,\omega_d).
$$
Then
$$
\int_{a_i}^{b_i}h_{\omega_{-i}}(t)\,dt=0.
$$
By Lemma~\ref{lem:half_oscillation},
$$
\sup_{s\in[a_i,b_i]}
\left|
\int_{a_i}^{s}h_{\omega_{-i}}(t)\,dt
\right|
\le
\frac12\int_{a_i}^{b_i}|h_{\omega_{-i}}(t)|\,dt.
$$
Thus
$$
\|\Psi_i\|_{L^\infty(S)}
\le
\frac{1}{2k}
\operatorname*{ess\,sup}_{\omega_{-i}\in S_{-i}}
\int_{a_i}^{b_i}
|\varphi(\omega_1,\ldots,\omega_{i-1},t,\omega_{i+1},\ldots,\omega_d)|\,dt.
$$
The factor $1/(2k)$ comes from averaging over the $k$ selected coordinate
directions and applying the half-oscillation estimate on each slice.
Taking the maximum over $i\in I$ gives
\begin{equation}
\label{eq:vector_field_sup_bound_tv}
\|\Psi\|_{L^\infty(S;\ell_\infty)}
\le
\frac{1}{2k}
\Lambda_{\mathrm{sl}}^I(\varphi).
\end{equation}

Since $g\in BV(S)$ and $\Psi\in\mathrm{Lip}(\overline S;\mathbb R^d)$, the
Gauss--Green identity yields
$$
\int_S g\varphi\,d\omega
=
\int_S g\,\operatorname{div}\Psi\,d\omega
=
\int_{\partial S}\operatorname{Tr}(g)(\Psi\cdot n)\,d\mathcal H^{d-1}
-
\int_S\Psi\cdot d(Dg).
$$
The boundary term vanishes. Since only coordinates in $I$ contribute,
$$
\left|\int_S g\varphi\,d\omega\right|
\le
\sum_{i\in I}\|\Psi_i\|_{L^\infty(S)}\,|D_i g|(S).
$$
Using
$$
\|\Psi_i\|_{L^\infty(S)}
\le
\|\Psi\|_{L^\infty(S;\ell_\infty)}
$$
and \eqref{eq:vector_field_sup_bound_tv}, we obtain
$$
|\mathbb E_f[\varphi]|
=
\left|\int_S g\varphi\,d\omega\right|
\le
\frac{1}{2k}
\Lambda_{\mathrm{sl}}^I(\varphi)
\TV_I(g;S)
=
\frac{1}{2k}
\Lambda_{\mathrm{sl}}^I(\varphi)
\TV_I(f;S).
$$
This proves \eqref{eq:k_direction_tv_bound}. Since
$\TV_I(f;S)\le \TV_\infty(f;S)$, \eqref{eq:k_direction_tv_bound_tvinfty}
follows.

Finally, if $\|\varphi\|_{L^\infty(S)}\le h$, then for every $i\in I$,
$$
\operatorname*{ess\,sup}_{\omega_{-i}\in S_{-i}}
\int_{a_i}^{b_i}
|\varphi(\omega_1,\ldots,\omega_{i-1},t,\omega_{i+1},\ldots,\omega_d)|\,dt
\le
(b_i-a_i)h.
$$
Therefore
$$
\Lambda_{\mathrm{sl}}^I(\varphi)
\le
\left(\max_{i\in I}(b_i-a_i)\right)h,
$$
and substitution into \eqref{eq:k_direction_tv_bound} gives
\eqref{eq:k_direction_tv_bound_sup}.
\end{proof}

\end{document}
